\documentclass[12pt,a4paper]{report}

% -------------------------
% Packages
% -------------------------
\usepackage[utf8]{inputenc}
\usepackage[T1]{fontenc}
\usepackage[french]{babel}

\usepackage{geometry}
\geometry{
    left=3cm,
    right=2.5cm,
    top=2.5cm,
    bottom=2.5cm
}

\usepackage{setspace}
\onehalfspacing

\usepackage{graphicx}
\usepackage{float}
\usepackage{caption}
\usepackage{subcaption}

\usepackage{amsmath, amssymb}
\usepackage{booktabs}
\usepackage{array}
\usepackage{longtable}

\usepackage{xcolor}
\usepackage{hyperref}
\usepackage{url}

\usepackage{fancyhdr}
\usepackage{titlesec}
\usepackage{enumitem}

% -------------------------
% Configuration
% -------------------------
\hypersetup{
    colorlinks=true,
    linkcolor=black,
    urlcolor=blue,
    citecolor=black,
    % pdftitle={Titre du rapport},
    % pdfauthor={Nom Prénom}
}

\pagestyle{fancy}
\fancyhf{}
\fancyhead[L]{Interprétabilité et alignement}
\fancyhead[R]{Pier-Luc Dupéré}
\fancyfoot[C]{\thepage}

\setlength{\parindent}{0pt}
\setlength{\parskip}{6pt}


% -------------------------
% Début du document
% -------------------------
\begin{document}

% -------------------------
% Page de garde
% -------------------------
\begin{titlepage}

    \begin{center}

        \includegraphics[width=0.50\textwidth]{img/logo-universite.png}
        
        \vspace{1cm}

        {\Large \textbf{Université du Québec à Rimouski}}\\[0.3cm]
        {\large Département de mathématiques, d'informatique et de génie}\\[0.3cm]
        {\large 8INF976 - Sujet spécial en intelligence artificielle}\\[1.5cm]

        \rule{\textwidth}{0.5pt}\\[0.4cm]

        {\Huge \textbf{Interprétabilité et alignement}}\\[0.4cm]
        {\Large Rapport de lecture}\\[0.4cm]

        \rule{\textwidth}{0.5pt}\\[2cm]

        \begin{flushleft}
            \textbf{Réalisé par :} Pier-Luc Dupéré \\[0.4cm]
            \textbf{Encadré par :} Mehdi Adda, PhD 
        \end{flushleft}

        \vfill
    \end{center}

\end{titlepage}

\tableofcontents
% \listoffigures
% \listoftables

% -------------------------
% Corps du rapport
% -------------------------
\clearpage
\pagenumbering{arabic}

\chapter{Introduction}
L'entraînement des grands modèles de langage (LLM) repose d'abord sur l'apprentissage 
auto-supervisé à partir de vastes corpus textuels. Cette étape leur permet d'apprendre des 
régularités linguistiques et d'acquérir des connaissances. Mais, cela ne garantit pas qu'ils 
produisent des réponses utiles, pertinentes et conformes aux attentes des utilisateurs. 
Sans entraînement complémentaire, ils peuvent notamment générer des contenus toxiques. 
Comme présenter des informations erronées ou mal suivre des consignes complexes.

Des méthodes d'alignement visent à orienter le comportement des modèles à partir de préférences 
humaines. Leur déploiement systématique permettrait de réduire significativement le taux de réponses 
toxiques. 


\chapter{Reinforcement Learning from Human Feedback (RLHF) [1] [2]}
L'apprentissage par renforcement à partir de rétroactions humaines (RLHF) est une technique utilisée 
pour intégrer des informations issues d'êtres humains dans les systèmes d'intelligence artificielle. 

Le processus fondamental du RLHF comprend trois étapes. Premièrement, il est nécessaire d'entraîner 
un modèle de langage capable de répondre aux questions des utilisateurs. Deuxièmement, il faut 
recueillir des données sur les préférences humaines afin d'entraîner un modèle de récompense qui 
reflète ces préférences. Finalement, le modèle de langage peut être optimisé à l'aide d'un algorithme 
d'apprentissage par renforcement. Cela ce fait en générant des réponses qui sont ensuite évaluées 
par le modèle de récompense. 

\begin{figure}[H]
    \centering
    \includegraphics[width=0.95\textwidth]{img/RLHF.jpg}
    \caption{Une représentation du processus RLHF initial en trois étapes, comprenant le SFT, 
            un modèle de récompense, puis une phase d'optimisation.}
    \label{fig:RLHF}
\end{figure}



\section{Entraîner un modèle de langage}

Le Supervised Fine-Tuning (SFT) consiste à prendre un modèle de langage de base pré-entraîné 
et à l'ajuster sur un jeu de données soigneusement sélectionné. Il est composé de paires 
instruction-réponse de haute qualité. Cet ensemble de données SFT contient généralement des 
milliers, voire des dizaines de milliers d'exemples illustrant le comportement attendu.
On y retrouve des réponses claires aux questions, des interventions utiles face à des requêtes 
ainsi que des refus appropriés lorsqu'il s'agit de demandes potentiellement dangereuses.

Le SFT transforme le modèle de base initialement conçu pour prédire le jeton suivant. Ce système qui 
est capable de suivre des instructions et de générer des réponses structurées. Toutefois, le SFT 
présente des limites. Le modèle assimile la forme que doivent prendre des réponses pertinentes. 
Il n'est pas toujours en mesure de produire systématiquement un contenu de haute qualité. De plus, il 
peut avoir de la difficulté à généraliser ses capacités au-delà des exemples spécifiques présents 
dans l'ensemble de données SFT.

\section{Entraînement d'un modèle de récompense}

La formation du modèle de récompense consiste à créer un modèle capable de prédire les préférences 
humaines. Des annotateurs humains se voient présenter un prompt ainsi que plusieurs réponses générées 
par un modèle. Par la suite, ils doivent indiquer laquelle est la meilleure. Ces étiquettes de 
comparaison servent à entraîner un modèle de récompense. Généralement le modèle de langage doté d'une 
sortie scalaire qui attribue un score numérique à chaque prompt qui donne une réponse. On peut ainsi 
prédire quelle réponse serait préférable pour un humain.
 
On utilise généralement entre des milliers et des centaines de milliers d'étiquettes de comparaison. 
Une fois le modèle de récompense entraîné, le système dispose donc d'un modèle de langage capable de 
générer des réponses et d'un modèle qui estime leur préférence. Le score de récompense peut alors 
servir à optimiser le modèle de langage à l'étape suivante du processus RLHF.


\section{Optimisation par apprentissage par renforcement}

Le modèle de récompense attribue un score à ces réponses et l'algorithme Proximal Policy Optimization (PPO) 
ajuste les poids du modèle de politique afin d'accroître la récompense attendue. Un élément crucial 
est la pénalité liée à la divergence, qui empêche le modèle de politique de trop s'éloigner 
du modèle SFT initial. 

Cette pénalité est souvent fondée sur la divergence de Kullback-Leibler. Elle est calculée à partir 
des distributions de probabilité des jetons. Cette pénalité permet de maintenir le modèle dans un 
cadre de génération linguistique cohérent et permet l'orientation des résultats qui maximisent la 
récompense.

Enfin, le processus peut être répété.Cette approche est parfois appelée RLHF en ligne itératif. 
Elle nécessite toutefois un dispositif permettant de recueillir régulièrement des retours humains 
sur les réponses produites.


\chapter{Direct Preference Optimization (DPO) [3] [4]}

L'innovation fondamentale de DPO réside dans son approche d'optimisation directe. Plutôt que 
d'entraîner un modèle de récompense distinct, DPO utilise une fonction de perte de type entropie 
croisée binaire pour mettre à jour directement les poids du modèle à partir de données de préférence. 
Ce processus simplifié rend l'entraînement plus stable et plus efficace, tout en permettant d'obtenir 
des résultats comparables, voire supérieurs, à ceux de la méthode traditionnelle RLHF.

La fonction de perte de DPO

\begin{figure}[H]
    \centering
    \includegraphics[width=0.95\textwidth]{img/DPO.jpg}
    \caption{La fonction de perte de DPO qui optimise directement la politique du le modèle de langage 
            tout en utilisant des données de préférence}
    \label{fig:DPO}
\end{figure}

De plus, DPO réduit la consommation de mémoire VRAM lors de l'entraînement et accélère les temps de 
convergence. 

Le DPO réduit ainsi significativement les obstacles à l'entraînement de modèles basés sur les 
préférences humaines. Par ailleurs, DPO ne nécessite pratiquement aucun réglage d'hyperparamètres.


\begin{figure}[H]
    \centering
    \includegraphics[width=0.95\textwidth]{img/pytorch.jpg}
    \caption{Le code PyTorch correspondant à la fonction de perte DPO}
    \label{fig:pytorch}
\end{figure}

\chapter{SAE [5]}
Les objectifs de ce rapport sont les suivants :

\begin{itemize}
    \item méthode de pointe pour l'entraînement d'auto-encodeurs parcimonieux.
    \item on souligne des lois d'échelle bien définies et on étend l'approche à un grand nombre de variables latentes.
    \item métriques permettant d’évaluer la qualité des variables latentes et constatons que, selon ces métriques, les auto-encodeurs parcimonieux de plus grande taille sont généralement plus performants.
\end{itemize}





\chapter{Red Teaming Language Models with Language Models [6]}




% -------------------------
% Bibliographie
% -------------------------
\begin{thebibliography}{9}

\bibitem{site}
Nathan Lambert,
\textit{Reinforcement Learning from Human Feedback},
\url{https://arxiv.org/abs/2504.12501},
consulté le 9 novembre 2026.

\bibitem{site}
Hugging Face,
\textit{Illustrating Reinforcement Learning from Human Feedback (RLHF) },
\url{https://huggingface.co/blog/rlhf},
consulté le 9 novembre 2026.

\bibitem{site}
Rafailov, R., Sharma, A., Mitchell, E., Ermon, S., Manning, C. D. et Finn, C.,
\textit{Direct Preference Optimization: Your Language Model is Secretly a Reward Model},
\url{https://arxiv.org/abs/2305.18290},
consulté le 10 novembre 2026.

\bibitem{site}
Rafailov, R., Sharma, A., Mitchell, E., Ermon, S., Manning, C. D. et Finn, C.,
\textit{Direct Preference Optimization: Your Language Model is Secretly a Reward Model},
\url{https://huggingface.co/learn/smol-course/unit2/2},
consulté le 10 novembre 2026.

\bibitem{site}
Gao, L., la Tour, T. D., Tillman, H., Goh, G., Troll, R., Radford, A. et al.,
\textit{Direct Preference Optimization: Your Language Model is Secretly a Reward Model},
\url{https://arxiv.org/abs/2406.04093},
consulté le 10 novembre 2026.

\bibitem{site}
Perez, E., Huang, S., Song, F., Cai, T., Ring, R., Aslanides, J. et al.,
\textit{Red Teaming Language Models with Language Models},
\url{https://arxiv.org/abs/2202.03286},
consulté le 10 novembre 2026.

\end{thebibliography}

% -------------------------
% Annexes
% -------------------------
% \appendix

% \chapter{Annexe A : Code source}

% Présentez ici des extraits de code, des tableaux supplémentaires ou des documents complémentaires.

% \chapter{Annexe B : Informations supplémentaires}

% Ajoutez ici les éléments qui ne sont pas indispensables à la compréhension du rapport principal.

\end{document}
